\section{零点密度网格与完整高度余项}\label{sec:density}

本节从 QRH$(7/8)$、CS v2 和 Rosser 零点计数输入推出
第 \ref{sec:trapezoid} 节所需的统一零点预算。
核心是两次正权离散化：先沿零点实部方向使用密度计数，
再沿高度方向使用递减核。每一次离散化的误差方向都是向上；
高度网格之外的贡献由一个明确的几何级数估计，不能省略。

\subsection{重数约定与 CS v2 的精确使用方式}

只对非平凡零点 $\rho=\beta+\ii\gamma$ 计数，并始终乘重数
$m_\rho$。令
\[
 N(T)=\sum_{0<\gamma\leq T}m_\rho,\qquad
 N(\sigma,T)=\sum_{\substack{0<\gamma\leq T\\\beta\geq\sigma}}m_\rho,
 \qquad H=3\cdot10^{12}.
\]
这里 $N(T)$ 与前文的 $N_+(T)$ 是同一计数的简写，
不是包含正、负高度的总和；$N$ 的第一参数省略时表示全部正高度零点。
闭条件 $\gamma\leq T$ 与 $\beta\geq\sigma$ 均与当前 Lean 输入一致。
若文献采用一个更大的零点集合，限制到本文集合只会减小计数。
共轭对称和重数保持允许把正、负高度的范数和合并为正高度的两倍。
非平凡实零点不存在，故没有 $\gamma=0$ 的剩余项。

对 $r=0,\ldots,11$，CS v2 输入 \cite{CSv2} 在
\[
 \frac{16+r}{32}\leq\sigma\leq\frac{17+r}{32},\qquad T\geq H
\]
上给出
\begin{equation}\label{dens:CS}
 N(\sigma,T)\leq
 B_{1,r}T^{p(\sigma)}(\log T)^{p(\sigma)+2}
 +B_{2,r}(\log T)^2+B_{3,r}\log T,
 \qquad p(\sigma)=\frac{3(1-\sigma)}{2-\sigma}.
\end{equation}
表 \ref{dens:constants} 的有限小数都按所写的精确有理数使用，
不是再次拟合所得的浮点常数。本文承认此条预印本输入；
本节不证明它的上游推导。

\begin{table}[htbp]
\centering
\caption{当前证明实际使用的 CS v2 前十二行常数。}\label{dens:constants}
\begin{tabular}{c|rrr@{\qquad}c|rrr}
$r$&$B_{1,r}$&$B_{2,r}$&$B_{3,r}$
 &$r$&$B_{1,r}$&$B_{2,r}$&$B_{3,r}$\\ \hline
0&5.069&9.461&167.8&6&12.11&7.704&136.7\\
1&6.046&9.168&162.6&7&13.66&7.411&131.5\\
2&7.077&8.875&157.4&8&15.35&7.118&126.3\\
3&8.185&8.582&152.2&9&17.21&6.826&121.1\\
4&9.387&8.290&147.0&10&19.31&6.533&115.9\\
5&10.70&7.997&141.9&11&21.77&6.240&110.7
\end{tabular}
\end{table}

式 \eqref{dens:CS} 中对数的指数是 $p(\sigma)+2$，
其恒等式形式为 $(7-5\sigma)/(2-\sigma)$。
在有限网格证书里不能把它任意换成旧版本的指数；
后文若使用 $\log^3T$，只能作为由
$p(\sigma)+2\leq3$ 和 $\log T\geq1$ 推出的较粗上界。
精确输入和这个代数等式分别见
\leanref{DensityInput}{CSv2} 和
\leanref{DensityInput}{cs_log_exponent_eq}。

Rosser 输入在本文所用高度上蕴含简便估计
\begin{equation}\label{dens:rosser-simple}
 N(\e^u)\leq\e^u u\qquad(\e^u\geq H,\ u\geq1).
\end{equation}
具体地，记源文闭矩形计数为 $N_R(T)$，则 $N(T)\leq N_R(T)$。
由 Rosser 上界、$\log T\geq1$、$\log\log T\leq\log T$ 和
$\pi\geq1$，丢弃负项后有
\[
 N(T)\leq\frac T2\log T+0.580\log T+2.463
       \leq\left(\frac T2+4\right)\log T
       \leq T\log T.
\]
最后一步只需 $T\geq8$，在本文高度域内自动满足。
这里未把 $N(T)\leq T\log T$ 作为第八个外部输入；
对应 \leanref{LiteratureInputs}{totalZeroCount_le_mul_log}。
原始零点计数来源为 \cite{Rosser1941,Dusart2016}。

\subsection{实部网格：从单个指数到密度计数}

设
\[
 \frac12=\sigma_0\leq\sigma_1\leq\cdots\leq\sigma_J=\frac78
\]
是有限网格。每个左端点 $\sigma_j$ 指定一个适用的 CS 行
$r(j)$，即 $\sigma_j$ 属于该行的闭区间。
对 $A\geq0$ 定义
\begin{equation}\label{dens:weights}
 W_j(A)=\e^{(\sigma_{j+1}-1)A}
             -\e^{(\sigma_j-1)A}\geq0.
\end{equation}

\begin{lemma}[向上的阶梯界]\label{dens:staircase}
若 $\beta\leq7/8$，则
\begin{equation}\label{dens:staircase-eq}
 \e^{(\beta-1)A}\leq
 \e^{-A/2}+
 \sum_{j=0}^{J-1}\mathbf{1}_{\{\sigma_j\leq\beta\}}W_j(A).
\end{equation}
\end{lemma}
\begin{proof}
若 $\beta<\sigma_0$，右边至少为 $\e^{-A/2}$，结论由单调性得到。
否则，取最后一个满足 $\sigma_j\leq\beta$ 的网格位置。
右边的增量和望远镜相消，等于一个不小于
$\e^{(\beta-1)A}$ 的右端点值。若 $\beta=\sigma_J$，
所有增量都被计入，右边恰等于 $\e^{(\sigma_J-1)A}$。
在内部网格端点使用闭条件可能多计一个增量，但仍为上界。
\end{proof}

定义正高度的带权累积计数
\[
 F_A(T)=\sum_{0<\gamma\leq T}m_\rho\e^{(\beta-1)A}.
\]
QRH$(7/8)$ 使引理 \ref{dens:staircase} 可用于每一个零点。
将其乘以重数并求和，得到
\begin{equation}\label{dens:weighted-count}
 F_A(T)\leq\e^{-A/2}N(T)+
              \sum_{j=0}^{J-1}W_j(A)N(\sigma_j,T).
\end{equation}
这里无需零点恰好落在某个半开网格单元内；
使用累积计数已处理所有边界重合。

把 $T=\e^u$ 代入，记 $p_j=p(\sigma_j)$，并定义
\begin{align}
 C_j(u)&=
 B_{1,r(j)}\e^{p_ju}u^{p_j+2}
       +B_{2,r(j)}u^2+B_{3,r(j)}u,\label{dens:C}\\
 \Phi_A(u)&=\e^{-A/2+u}u+
                     \sum_{j=0}^{J-1}W_j(A)C_j(u).
                                                   \label{dens:Phi}
\end{align}
当 $\e^u\geq H$、$u\geq1$ 时，
\begin{equation}\label{dens:F-Phi}
 F_A(\e^u)\leq\Phi_A(u).
\end{equation}
至此，所有零点位置都已被一个显式的、有限项的实函数替换。
对应 \leanref{FiniteDensityGrid}{exp_le_grid_staircase}、
\leanref{FiniteDensityGrid}{positiveZeroWeight_le_density_grid}
和 \leanref{FiniteDensityGrid}{positiveZeroWeight_le_gridPhi}。

\subsection{密度上界沿高度方向的增长率}

为了控制无限余项，先证明一个不依赖任何数值采样的增长界。
\begin{lemma}\label{dens:growth}
若 $U\geq6$、$v\geq0$，则
\begin{equation}\label{dens:growth-eq}
 \Phi_A(U+v)\leq\e^{3v/2}\Phi_A(U).
\end{equation}
\end{lemma}
\begin{proof}
对 $u>0$、$v\geq0$、$r\geq0$，由
$\log(1+v/u)\leq v/u$ 得
\[
 \e^{p(u+v)}(u+v)^r
 \leq \e^{(p+r/u)v}\e^{pu}u^r.
\]
在 $1/2\leq\sigma_j\leq7/8$ 上，$0\leq p_j\leq1$，
且 $2\leq p_j+2\leq3$。
因此 $p_j+(p_j+2)/U\leq1+3/6=3/2$。
同样的估计也适用于 $u^2$、$u$ 和基线 $\e^u u$：
相应的 $(p,r)$ 分别为 $(0,2)$、$(0,1)$、$(1,1)$。
式 \eqref{dens:Phi} 中的系数均非负，逐项相加即得结论。
\end{proof}
特别地，将高度乘以 $4$ 时，对数高度增加 $\log4$，
从而 $\Phi_A(U+\log4)\leq8\Phi_A(U)$。
该界对应 \leanref{DensityGridGrowth}{densityGridPhi_shift_le}。

\subsection{高度网格及其不能遗漏的剩余权}

选择单调对数高度网格
$u_0\leq u_1\leq\cdots\leq u_M=U$，并要求
\begin{equation}\label{dens:height-conditions}
 \e^{u_0}\geq H,\qquad u_0\geq1,\qquad U\geq6,\qquad
 g_q(\e^{u_0})=1.
\end{equation}
记 $t_i=\e^{u_i}$、$R=\e^U$，并定义非负质量
\begin{equation}\label{dens:mu}
 \mu_i=g_q(t_i)-g_q(t_{i+1})\geq0,\qquad 0\leq i<M.
\end{equation}
对于任意 $t>0$，有
\begin{equation}\label{dens:height-staircase}
 g_q(t)\leq
 \sum_{i=0}^{M-1}\mu_i\mathbf{1}_{\{t\leq t_{i+1}\}}
        +g_q(\max\{t,R\}).
\end{equation}
若 $t\leq t_0$，右边望远镜相消为 $1$；
若 $t_i<t\leq t_{i+1}$，右边不小于 $g_q(t_i)$；
若 $t>R$，右边等于 $g_q(t)$。
这证明了 \eqref{dens:height-staircase}，并说明余项具有明确的含义：
\emph{即使一个零点位于 $R$ 以下，它在有限高度和之外仍留下
$g_q(R)$ 这一份权。}
因此，仅估计 $\gamma>R$ 的零点尾和不足以完成本证明。

对任意有限截断 $T$，将 \eqref{dens:height-staircase}
乘以 $m_\rho\e^{(\beta-1)A}$ 并对正高度零点求和，
有限网格部分不超过
\[
 \sum_{i=0}^{M-1}\mu_i F_A(t_{i+1})
 \leq\sum_{i=0}^{M-1}\mu_i\Phi_A(u_{i+1}).
\]
剩余部分使用
\begin{equation}\label{dens:quadratic-residual}
 g_q(\max\{t,R\})\leq
             \frac{c_2q^2}{\max\{t,R\}^2}.
\end{equation}

\begin{lemma}[完整二次余项]\label{dens:residual}
在 \eqref{dens:height-conditions} 的条件下，对每个 $T$，
\begin{equation}\label{dens:residual-eq}
 \sum_{0<\gamma\leq T}
 \frac{m_\rho\e^{(\beta-1)A}c_2q^2}{\max\{\gamma,R\}^2}
 \leq16c_2q^2\e^{-2U}\Phi_A(U).
\end{equation}
\end{lemma}
\begin{proof}
若截断集合为空，结论显然。否则选取整数 $N\geq0$，
使集合内所有高度均不超过 $R4^{N+1}$。
对于这些高度，逐项有
\[
 \frac{c_2q^2}{\max\{t,R\}^2}
 \leq\frac{c_2q^2}{R^2}
       \sum_{k=0}^N16^{-k}\mathbf{1}_{\{t\leq R4^{k+1}\}}.
\]
若 $t\leq R$，右边的 $k=0$ 项已足够；
若 $R4^k<t\leq R4^{k+1}$，右边第 $k$ 项已足够。
对零点求和，并用 \eqref{dens:F-Phi} 与引理
\ref{dens:growth}，得到
\begin{align*}
 \text{左边}
 &\leq \frac{c_2q^2}{R^2}
   \sum_{k=0}^N16^{-k}\Phi_A(U+(k+1)\log4)\\
 &\leq \frac{c_2q^2}{R^2}\Phi_A(U)
   \sum_{k=0}^N16^{-k}8^{k+1}\\
 &\leq \frac{c_2q^2}{R^2}\Phi_A(U)\,
          8\sum_{k=0}^{\infty}2^{-k}
   =16c_2q^2\e^{-2U}\Phi_A(U).
\end{align*}
得到的常数与 $T$、$N$ 均无关。
\end{proof}
这给出常数 $16$ 的完整来源。它已经包含高度 $R$ 以下的剩余权，
以及 $R$ 以上全部高度；后续不会再另留一个未定义的尾项。
对应 \leanref{GeometricHeightTail}{finite_quadratic_residual_bound}
和 \leanref{TrapezoidHeightBound}{positive_zero_quadratic_residual_le}。

\subsection{从密度预算回到梯形零点和}

定义可供数值认证的完整正高度预算
\begin{equation}\label{dens:B}
 \mathcal B(A,q;\boldsymbol u)=
 \sum_{i=0}^{M-1}
   \bigl[g_q(\e^{u_i})-g_q(\e^{u_{i+1}})\bigr]\Phi_A(u_{i+1})
       +16c_2q^2\e^{-2U}\Phi_A(U).
\end{equation}
前面的论证表明，对任意 $T$，
\begin{equation}\label{dens:positive-final}
 \sum_{0<\gamma\leq T}m_\rho\e^{(\beta-1)A}g_q(\gamma)
 \leq\mathcal B(A,q;\boldsymbol u).
\end{equation}

\begin{proposition}[完整零点预算]\label{dens:zero-budget}
设 QRH$(7/8)$、CS v2 和 Rosser 输入成立。
若 $x=\e^L$、$0\leq A\leq L$、$h>0$、
$q\geq(x+h)/h$，网格满足以上条件，则对任意 $T$，
\begin{equation}\label{dens:zero-budget-eq}
 |Z_T(x,h)|\leq2h\,\mathcal B(A,q;\boldsymbol u).
\end{equation}
所以第 \ref{sec:trapezoid} 节的判据可取
$Z=2\mathcal B(A,q;\boldsymbol u)$。
\end{proposition}
\begin{proof}
由 QRH，$\beta\leq7/8<1$，故 $L\geq A$ 蕴含
$\e^{(\beta-1)L}\leq\e^{(\beta-1)A}$。
用引理 \ref{tr:three-kernels} 的合并界、三角不等式和共轭配对，
\[
 |Z_T(x,h)|
 \leq2h\sum_{0<\gamma\leq T}
                  m_\rho\e^{(\beta-1)A}g_q(\gamma).
\]
再用 \eqref{dens:positive-final} 即可。
\end{proof}
对应 \leanref{TrapezoidHeightBound}{positive_zero_envelope_le_heightBudget}
和
\leanref{TrapezoidHeightBound}{norm_truncatedTrapezoidZeroSum_le_heightBudget}。
结合命题 \ref{tr:prime-criterion}，每个待认证单元的充分条件现为
\[
 2\mathcal B(A,q;\boldsymbol u)+P_k(A,C)+V_A<1.
\]
下一节把这里的有限实数表达式转为可审阅、可逐项重放的有理证书。
